\section{Equality cases}
\label{sec:equality-cases}

We now determine all graphs attaining equality in \eqref{eq:main-bound}.
The argument uses two additional exact identities for the matrices
$V_\sigma$, followed by a count of common neighbors.

\begin{lemma}\label{lem:equality-flag-identities}
Let $k\in\{1,2\}$, $\ell=\lfloor(8+k)/2\rfloor$, and $\sigma=([k],\varnothing)$.
Every flag $P\in\mathcal F_{\sigma,\ell}$ has three unlabeled vertices.
Define
\[
d(P)=\big|\{v\in V(P)\setminus[k]:\{1,v\}\in E(P)\}\big|.
\]
For $k=2$, also define
\[
c(P)=\big|\{v\in V(P)\setminus[2]:
\{1,v\}\in E(P),\ \{2,v\}\in E(P)\}\big|.
\]
For the fixed matrices $V_\sigma$ used in
\eqref{eq:stored-matrix-construction}, the following statements hold.
If $\sigma=([1],\varnothing)$, there exists a rational column vector
$b\in\mathbb Q^{10}$ such that
\begin{equation}\label{eq:equality-one-root-vector}
V_\sigma b
=\big(25d(P)(d(P)-1)-11d(P)\big)_{P\in\mathcal F_{\sigma,\ell}}.
\end{equation}
If $\sigma=([2],\varnothing)$, there exists a rational column vector
$b\in\mathbb Q^{54}$ such that
\begin{equation}\label{eq:equality-two-root-vector}
V_\sigma b
=\big(25c(P)(c(P)-1)-14c(P)+3d(P)\big)_{P\in\mathcal F_{\sigma,\ell}}.
\end{equation}
\end{lemma}
The proof of Lemma~\ref{lem:equality-flag-identities} is entirely computational;
see \href{https://github.com/Diudiutse/triangle-free-signless-laplacian/blob/main/programs/verify_equality.py}{\texttt{programs/verify\_equality.py}}.

We also use the following uniqueness theorem.
\begin{lemma}[{\cite[Theorem~6.4]{Gewirtz1969}}]
\label{lem:higman-sims-uniqueness}
Up to isomorphism, the Higman--Sims graph is the unique triangle-free
$22$-regular graph on $100$ vertices in which any two distinct
nonadjacent vertices have exactly $6$ common neighbors.
\end{lemma}
In the original notation, this graph is $r(2,6,22,100)$.
The parameters specify diameter $2$, exactly $6$ paths of length $2$
between vertices at distance $2$, degree $22$, and order $100$.
The conditions in Lemma~\ref{lem:higman-sims-uniqueness} imply diameter
$2$; they also imply girth $4$, since two nonadjacent vertices and two
of their common neighbors form a $4$-cycle, while triangles are excluded.

\noindent\emph{Proof of the equality statement in Theorem~\ref{thm:bound}.}
As shown in \cite{deLimaNikiforovOliveira}, for every integer
$m\geq1$, the graph $\mathrm{HS}^m$ is triangle-free and satisfies
\[
|V(\mathrm{HS}^m)|=100m,
\qquad \qmin(\mathrm{HS}^m)=14m=\frac7{50}|V(\mathrm{HS}^m)|.
\]
Graph isomorphisms preserve the signless Laplacian spectrum, so the
same equality holds for every $G\cong \mathrm{HS}^m$.

Conversely, suppose that $\qmin(G)=7n/50$.
Write $V=V(G)$, $A=A_G$, and
\[
Q_G=\operatorname{diag}_{v\in V}(\deg_G(v))+A.
\]
Every eigenvalue of $Q_G/n-7I/50$ is nonnegative.
By \eqref{eq:new-F-matrix}, for every real column vector $u\in\R^5$,
\begin{align*}
u^{\mathsf T}\mathsf F u
&=\E_{x,y}\left[
\frac{A_{xy}}n
(U^{xy}u)^{\mathsf T}
\left(\frac{Q_G}n-\frac7{50}I\right)
(U^{xy}u)\right]\geq0.
\end{align*}
Thus $\mathsf F$ is positive semidefinite.
Let $B$ be as in the proof of \eqref{eq:main-bound}.
Hence $\tr(B\mathsf F)\geq0$.
By the proof of \eqref{eq:main-bound}, $\tr(B\mathsf F)\leq0$.
Therefore
\begin{equation}\label{eq:equality-zero-trace}
\tr(B\mathsf F)=0.
\end{equation}

For each type $\sigma\in\mathcal T_k$, $0\leq k\leq6$, set
$\ell=\lfloor(8+k)/2\rfloor$ and define the column vector
\[
p_\sigma(\boldsymbol r)
=\big(\Pr(P\mid\boldsymbol r)\big)_{P\in\mathcal F_{\sigma,\ell}}
\qquad(\boldsymbol r\in V^k).
\]
In the following expectation, $\boldsymbol r$ is uniformly distributed on $V^k$.
\begin{samepage}
By \eqref{eq:equality-zero-trace}, \eqref{eq:trace-flag-decomposition},
\eqref{eq:flag-entry-average}, and \eqref{eq:stored-matrix-construction},
we obtain
\begin{align}
0&=\tr(B\mathsf F)\notag\\
&=-\sum_{k=0}^6\sum_{\sigma\in\mathcal T_k}
\tr\left(M_\sigma\sum_{H\in\mathcal T_8}
\Pr(G_{\boldsymbol X}\cong H)C_H^\sigma\right)\notag\\
&=-8!\sum_{k=0}^6\sum_{\sigma\in\mathcal T_k}
\tr\Bigl(M_\sigma\E_{\boldsymbol r}\Bigl[
\ind_{G_{\boldsymbol r}=\sigma}
\bigl(\Pr(P\mid\boldsymbol r)\Pr(P'\mid\boldsymbol r)\bigr)
_{P,P'\in\mathcal F_{\sigma,\ell}}\Bigr]\Bigr)\notag\\
&=-\sum_{k=0}^6\sum_{\sigma\in\mathcal T_k}
\frac{2^{80}}{a_\sigma^2}
\tr\Bigl(V_\sigma(qI+X_\sigma)V_\sigma^{\mathsf T}
\E_{\boldsymbol r}\Bigl[\ind_{G_{\boldsymbol r}=\sigma}
\bigl(\Pr(P\mid\boldsymbol r)\Pr(P'\mid\boldsymbol r)\bigr)
_{P,P'\in\mathcal F_{\sigma,\ell}}\Bigr]\Bigr)\notag\\
&=-\sum_{k=0}^6\sum_{\sigma\in\mathcal T_k}
\frac{2^{80}}{a_\sigma^2}
\tr\Bigl(V_\sigma(qI+X_\sigma)V_\sigma^{\mathsf T}
\E_{\boldsymbol r}\Bigl[\ind_{G_{\boldsymbol r}=\sigma}
p_\sigma(\boldsymbol r)p_\sigma(\boldsymbol r)^{\mathsf T}\Bigr]\Bigr)\notag\\
&=-\sum_{k=0}^6\sum_{\sigma\in\mathcal T_k}
\frac{2^{80}}{a_\sigma^2}
\E_{\boldsymbol r}\Bigl[\ind_{G_{\boldsymbol r}=\sigma}
\tr\Bigl(V_\sigma(qI+X_\sigma)V_\sigma^{\mathsf T}
p_\sigma(\boldsymbol r)p_\sigma(\boldsymbol r)^{\mathsf T}\Bigr)\Bigr]\notag\\
&=-\sum_{k=0}^6\sum_{\sigma\in\mathcal T_k}
\frac{2^{80}}{a_\sigma^2}
\E_{\boldsymbol r}\Bigl[\ind_{G_{\boldsymbol r}=\sigma}
p_\sigma(\boldsymbol r)^{\mathsf T}
V_\sigma(qI+X_\sigma)V_\sigma^{\mathsf T}
p_\sigma(\boldsymbol r)\Bigr].
\label{eq:equality-flag-zero-sum}
\end{align}
\end{samepage}
Every $qI+X_\sigma$ is positive definite by
condition~(\ref{cond:integer-flag-positive}).
Since all summands inside these finite averages are nonnegative and
each tuple has probability $n^{-k}>0$, it follows that for every
$\boldsymbol r\in V^k$ with $G_{\boldsymbol r}=\sigma$,
\begin{equation}\label{eq:equality-root-kernel}
V_\sigma^{\mathsf T}p_\sigma(\boldsymbol r)=0.
\end{equation}
Consequently, for every $\boldsymbol r\in V^k$ with
$G_{\boldsymbol r}=\sigma$ and every column vector $w$ of the appropriate
dimension,
\begin{equation}\label{eq:equality-column-consequence}
\sum_{P\in\mathcal F_{\sigma,\ell}}
(V_\sigma w)_P\Pr(P\mid\boldsymbol r)
=w^{\mathsf T}V_\sigma^{\mathsf T}p_\sigma(\boldsymbol r)=0.
\end{equation}

Fix $x\in V$ and take $\sigma=([1],\varnothing)$.
For $P\in\mathcal F_{\sigma,4}$, we write $\Pr(P\mid x)$ for
$\Pr(P\mid(x))$ in \eqref{eq:conditional-probability}, viewing $(x)$
as a one-dimensional root vector.
Let $\boldsymbol Y=(Y_2,Y_3,Y_4)$ be uniformly distributed on $V^3$.
Then
\begingroup
\allowdisplaybreaks[1]
\begin{align}
&\sum_{P\in\mathcal F_{\sigma,4}}d(P)\Pr(P\mid x)\notag\\
&=\sum_{P\in\mathcal F_{\sigma,4}}d(P)
\frac{\big|\{\boldsymbol Y\in V^3:
(G_{(x,\boldsymbol Y)},\iota)\cong_\sigma P\}\big|}{n^3}\notag\\
&=\sum_{P\in\mathcal F_{\sigma,4}}d(P)
\E_{\boldsymbol Y}\big[\ind_{\{(G_{(x,\boldsymbol Y)},\iota)\cong_\sigma P\}}\big]\notag\\
&=\E_{\boldsymbol Y}\big[d(G_{(x,\boldsymbol Y)})\big]\notag\\
&=\E_{\boldsymbol Y}\left[\sum_{i=2}^4 A_{xY_i}\right]\notag\\
&=\sum_{i=2}^4\E_{Y_i}A_{xY_i}
=\frac3n\sum_{z\in V}A_{xz}
=\frac{3\deg_G(x)}n.
\label{eq:equality-degree-first-moment}
\end{align}
Using $A_{xY_i}^2=A_{xY_i}$, we also have
\begin{align}
&\sum_{P\in\mathcal F_{\sigma,4}}d(P)(d(P)-1)\Pr(P\mid x)\notag\\
&=\sum_{P\in\mathcal F_{\sigma,4}}d(P)(d(P)-1)
\frac{\big|\{\boldsymbol Y\in V^3:
(G_{(x,\boldsymbol Y)},\iota)\cong_\sigma P\}\big|}{n^3}\notag\\
&=\sum_{P\in\mathcal F_{\sigma,4}}d(P)(d(P)-1)
\E_{\boldsymbol Y}\big[\ind_{\{(G_{(x,\boldsymbol Y)},\iota)\cong_\sigma P\}}\big]\notag\\
&=\E_{\boldsymbol Y}\big[
 d(G_{(x,\boldsymbol Y)})(d(G_{(x,\boldsymbol Y)})-1)\big]\notag\\
&=\E_{\boldsymbol Y}\left[
\left(\sum_{i=2}^4 A_{xY_i}\right)^2-\sum_{i=2}^4 A_{xY_i}\right]\notag\\
&=\sum_{\substack{2\leq i,j\leq4\\i\ne j}}
\E_{Y_i,Y_j}[A_{xY_i}A_{xY_j}]\notag\\
&=\sum_{\substack{2\leq i,j\leq4\\i\ne j}}
\left(\frac1n\sum_{z\in V}A_{xz}\right)
\left(\frac1n\sum_{z'\in V}A_{xz'}\right)
=6\left(\frac{\deg_G(x)}n\right)^2.
\label{eq:equality-degree-second-moment}
\end{align}
Taking $w=b$ in \eqref{eq:equality-column-consequence}, and using
\eqref{eq:equality-one-root-vector},
\eqref{eq:equality-degree-first-moment}, and
\eqref{eq:equality-degree-second-moment}, we obtain
\begin{align*}
0&=\sum_{P\in\mathcal F_{\sigma,4}}(V_\sigma b)_P\Pr(P\mid x)\\
&=\sum_{P\in\mathcal F_{\sigma,4}}
\big(25d(P)(d(P)-1)-11d(P)\big)\Pr(P\mid x)\\
&=25\cdot6\left(\frac{\deg_G(x)}n\right)^2
-11\cdot\frac{3\deg_G(x)}n\\
&=150\frac{\deg_G(x)}n
\left(\frac{\deg_G(x)}n-\frac{11}{50}\right).
\end{align*}
The equality $\qmin(G)=7n/50>0$ excludes isolated vertices,
so $\deg_G(x)>0$.
It follows that $G$ is regular, with
\begin{equation}\label{eq:equality-regularity}
\deg_G(x)=\frac{11n}{50}\qquad(x\in V).
\end{equation}

We next determine the common-neighbor counts.
Write $N_G(x)=\{z\in V:A_{xz}=1\}$.
For $x,y\in V$ with $A_{xy}=0$, take
$\sigma=([2],\varnothing)$.
Here $\Pr(P\mid x,y)$ denotes $\Pr(P\mid(x,y))$ in
\eqref{eq:conditional-probability}, with root vector $(x,y)$.
For $\boldsymbol Y=(Y_3,Y_4,Y_5)$ uniformly distributed on $V^3$,
\begin{align*}
&\sum_{P\in\mathcal F_{\sigma,5}}c(P)\Pr(P\mid x,y)\\
&=\sum_{P\in\mathcal F_{\sigma,5}}c(P)
\frac{\big|\{\boldsymbol Y\in V^3:
(G_{(x,y,\boldsymbol Y)},\iota)\cong_\sigma P\}\big|}{n^3}\\
&=\sum_{P\in\mathcal F_{\sigma,5}}c(P)
\E_{\boldsymbol Y}\big[\ind_{\{(G_{(x,y,\boldsymbol Y)},\iota)\cong_\sigma P\}}\big]\\
&=\E_{\boldsymbol Y}\big[c(G_{(x,y,\boldsymbol Y)})\big]\\
&=\E_{\boldsymbol Y}\left[\sum_{i=3}^5 A_{xY_i}A_{yY_i}\right]\\
&=\sum_{i=3}^5\E_{Y_i}[A_{xY_i}A_{yY_i}]
=\frac3n\sum_{z\in V}A_{xz}A_{yz}\\
&=\frac{3|N_G(x)\cap N_G(y)|}{n}.
\end{align*}
Similarly,
\begin{align*}
&\sum_{P\in\mathcal F_{\sigma,5}}c(P)(c(P)-1)\Pr(P\mid x,y)\\
&=\sum_{P\in\mathcal F_{\sigma,5}}c(P)(c(P)-1)
\frac{\big|\{\boldsymbol Y\in V^3:
(G_{(x,y,\boldsymbol Y)},\iota)\cong_\sigma P\}\big|}{n^3}\\
&=\sum_{P\in\mathcal F_{\sigma,5}}c(P)(c(P)-1)
\E_{\boldsymbol Y}\big[\ind_{\{(G_{(x,y,\boldsymbol Y)},\iota)\cong_\sigma P\}}\big]\\
&=\E_{\boldsymbol Y}\big[
 c(G_{(x,y,\boldsymbol Y)})(c(G_{(x,y,\boldsymbol Y)})-1)\big]\\
&=\E_{\boldsymbol Y}\left[
\left(\sum_{i=3}^5 A_{xY_i}A_{yY_i}\right)^2
-\sum_{i=3}^5 A_{xY_i}A_{yY_i}\right]\\
&=\sum_{\substack{3\leq i,j\leq5\\i\ne j}}
\E_{Y_i,Y_j}[A_{xY_i}A_{yY_i}A_{xY_j}A_{yY_j}]\\
&=\sum_{\substack{3\leq i,j\leq5\\i\ne j}}
\left(\frac1n\sum_{z\in V}A_{xz}A_{yz}\right)
\left(\frac1n\sum_{z'\in V}A_{xz'}A_{yz'}\right)
\\
&=6\left(\frac{|N_G(x)\cap N_G(y)|}{n}\right)^2.
\end{align*}
For the degree count in the same flags,
\begin{align*}
&\sum_{P\in\mathcal F_{\sigma,5}}d(P)\Pr(P\mid x,y)\\
&=\sum_{P\in\mathcal F_{\sigma,5}}d(P)
\frac{\big|\{\boldsymbol Y\in V^3:
(G_{(x,y,\boldsymbol Y)},\iota)\cong_\sigma P\}\big|}{n^3}\\
&=\sum_{P\in\mathcal F_{\sigma,5}}d(P)
\E_{\boldsymbol Y}\big[\ind_{\{(G_{(x,y,\boldsymbol Y)},\iota)\cong_\sigma P\}}\big]\\
&=\E_{\boldsymbol Y}\big[d(G_{(x,y,\boldsymbol Y)})\big]\\
&=\E_{\boldsymbol Y}\left[\sum_{i=3}^5 A_{xY_i}\right]\\
&=\sum_{i=3}^5\E_{Y_i}A_{xY_i}
=\frac3n\sum_{z\in V}A_{xz}
=\frac{3\deg_G(x)}n=\frac{33}{50}.
\end{align*}
Using \eqref{eq:equality-two-root-vector} in
\eqref{eq:equality-column-consequence}, we therefore obtain
\begin{align*}
0&=\sum_{P\in\mathcal F_{\sigma,5}}
\big(25c(P)(c(P)-1)-14c(P)+3d(P)\big)\Pr(P\mid x,y)\\
&=25\cdot6\left(\frac{|N_G(x)\cap N_G(y)|}{n}\right)^2-14\cdot3\frac{|N_G(x)\cap N_G(y)|}{n}+3\cdot\frac{33}{50}\\
 &=150\left(\frac{|N_G(x)\cap N_G(y)|}{n}\right)^2-42\frac{|N_G(x)\cap N_G(y)|}{n}+\frac{99}{50}\\
 &=150\left(\frac{|N_G(x)\cap N_G(y)|}{n}-\frac3{50}\right)
       \left(\frac{|N_G(x)\cap N_G(y)|}{n}-\frac{11}{50}\right).
\end{align*}
\endgroup
Hence
\begin{equation}\label{eq:equality-common-neighbors}
|N_G(x)\cap N_G(y)|\in
\left\{\frac{3n}{50},\frac{11n}{50}\right\}
\qquad(A_{xy}=0).
\end{equation}
For adjacent $x,y$, the common-neighbor count is zero by
triangle-freeness.

To recover the structure of $G$, partition $V$ into the classes
\[
W_x=\{y\in V:N_G(y)=N_G(x)\}.
\]
Each class is independent, since adjacent vertices cannot have equal
neighbors.
Moreover, for $x'\in W_x$ and $y'\in W_y$,
\[
A_{x'y'}=A_{xy'}=A_{xy},
\]
so two distinct classes induce either a complete bipartite graph or
no edges.
Put $d=11n/50$.
By regularity, all neighborhoods have size $d$, and thus
\[
|N_G(x)\cap N_G(y)|=d
\quad\Longleftrightarrow\quad y\in W_x.
\]
For $x\in V$,
\begin{align}
d^2
&=\sum_{z\in N_G(x)}\deg_G(z)
 =\sum_{z\in N_G(x)}\sum_{y\in V}A_{zy}\notag\\
&=\sum_{y\in V}\sum_{z\in N_G(x)}A_{zy}
 =\sum_{y\in V}|N_G(x)\cap N_G(y)|\notag\\
&=\sum_{y\in W_x}d+\sum_{y\in N_G(x)}0
 +\sum_{y\in V\setminus(W_x\cup N_G(x))}\frac{3n}{50}\notag\\
&=|W_x|d+(n-|W_x|-d)\frac{3n}{50}\notag\\
&=\frac{8n}{50}|W_x|+\frac{117n^2}{2500}.
\label{eq:equality-common-neighbor-total}
\end{align}
Hence
\begin{equation}\label{eq:equality-class-size}
|W_x|=\frac{d^2-117n^2/2500}{8n/50}
=\frac{4n^2/2500}{8n/50}=\frac n{100}.
\end{equation}
Thus $\frac n{100}$ is a positive integer and there are exactly $100$ classes.

Let $R$ be the graph whose vertices are these classes, with two
classes adjacent if all edges between them are present in $G$.
The classes are independent and all have size $\frac n{100}$, so $G\cong R^{n/100}$.
The graph $R$ is triangle-free, since any triangle in $R$ would lift
to a triangle in $G$, and each vertex of $R$ has degree
\[
\frac{d}{n/100}=\frac{11n/50}{n/100}=22.
\]
If $W_x,W_y$ are distinct nonadjacent vertices of $R$, then $x,y$
are nonadjacent in $G$ and have different neighborhoods.
By \eqref{eq:equality-common-neighbors}, they have exactly $3n/50$
common neighbors in $G$.
These common neighbors form a disjoint union of classes of size $\frac n{100}$,
so
\[
|N_R(W_x)\cap N_R(W_y)|
=\frac{|N_G(x)\cap N_G(y)|}{n/100}
=\frac{3n/50}{n/100}=6.
\]
Lemma~\ref{lem:higman-sims-uniqueness} now gives $R\cong \mathrm{HS}$.
Therefore $G\cong R^{n/100}\cong \mathrm{HS}^{n/100}$, completing the proof.\hfill$\square$
